\section{Related Work}
\label{sec:related_work}

The literature on parameter merging, Mixture-of-Experts (MoE) routing, and modular serving has progressed through several distinct evolutionary stages. In this section, we audit this progression and position our proposed framework within the historical landscape of multi-task model merging and stream serving.

\subsection{Static and Quantization-Aware Model Merging}
Model merging seeks to combine multiple specialized deep networks into a single multi-task model without additional training. Early foundational work focused on weight-space alignments, introducing weight-space diffeomorphisms (FoldMerge \cite{foldmerge_2024}) and polynomial splines on optimization landscapes (PolyMerge \cite{polymerge_2024}). As edge deployment became prominent, quantization-aware merging methods (Q-Merge \cite{qmerge_2024}) emerged to preserve performance under low-bit precisions, though independent audits soon exposed significant robustness vulnerabilities under extreme quantization \cite{qmerge_audit_2025}. These static merging methods rely on standard average-pooling or task arithmetic \cite{task_arithmetic_ilharco, wortsman2022model, merge_ties, merge_dare}, which cannot dynamically adapt parameter ensembling coefficients to heterogeneous query streams at runtime.

\subsection{Sparsity and Regularization in Model Merging}
To mitigate representation redundancy and task suite biases in joint parameter ensembling, researchers introduced sparsity and regularization. Sparse Task Arithmetic \cite{sparse_task_arithmetic_2025} and SuiteMerge \cite{suitemerge_2025} addressed interference by isolating orthogonal parameter subspaces. To prevent out-of-distribution overfitting, Task-Space Anchor Regularization (TSAR) \cite{tsar_2025} leveraged stable early-layer representations to anchor weight combinations. While regularized merging methods improve robustness, their static nature limits their utility when processing highly dynamic, time-varying query streams.

\subsection{Stateless Edge Serving and Dynamic Routing}
To enable single-pass model adaptation under sequential edge workloads, systems shifted toward dynamic, input-conditioned routing. SABLE \cite{sable_2025} and SPS-ZCA \cite{sps_zca_2025} proposed stateless dynamic routing using early-layer activation centroids to compute layer-wise ensembling weights in a single forward pass under $O(1)$ latency. However, because these methods compute routing weights independently at each layer, they suffer from the \emph{routing jitter paradox} \cite{chemmerge_2026}: noisy query activations cause the ensembling coefficients to oscillate violently across adjacent layers. This high-frequency layer-wise jitter triggers cascading representation collapse, severely degrading multi-task serving quality.

\subsection{Stateful Kinetics and the Hysteresis Trade-off}
To resolve the routing jitter paradox, stateful serv-time models were proposed. Inspired by biochemistry, ChemMerge \cite{chemmerge_2026} modeled active adapter ensembling using continuous-time non-equilibrium chemical kinetics (employing Arrhenius-like reaction rates and ODEs) to filter high-frequency oscillations. Momentum-Merge \cite{momentum_merge_2026} deconstructed this biochemical complexity, proving its equivalence to a simple Exponential Moving Average (EMA) applied across sequential samples. PAC-Kinetics \cite{pac_kinetics_2026} subsequently established a rigorous learning-theoretic and control-theoretic foundation for stateful kinetics, employing Lyapunov stability and a Catoni-type PAC-Bayesian bound for stationary $\beta$-mixing stochastic processes. 

However, by maintaining an internal state across sequence samples, these kinetics-based methods apply a low-pass filter in the \emph{temporal} (sample) domain to achieve smoothness in the \emph{spatial} (layer) domain. Consequently, when deployed under rapid task switches, they suffer from severe representational hysteresis and inertial serving lag, leading to major accuracy drops.

\subsection{Physics-Inspired Machine Learning and Path Integrals}
Our work draws on the rich history of mapping physical systems to probabilistic graphical models. The formulation of network depth as a 1D lattice shares structural mathematical properties with the 1D Ising model in statistical mechanics \cite{ising1925beitrag} and discrete Euclidean path integrals \cite{feynman1948space, feynman2010quantum}. While path integrals have been used in machine learning for continuous trajectory optimization or diffusion modeling, QPathMerge is the first framework to model deep network ensembling as a path integral over network depth. By solving this path integral exactly via sum-product message passing (Belief Propagation) \cite{belief_prop_pearl, exact_sum_product}, we decouple spatial layer smoothing from temporal sample tracking. This enables a zero-lag, hysteresis-free serving controller that delivers near-oracle smoothness across network layers.
